% GENERATED FILE. Edit canonical lessons, then run npm run generate:books.
% canonical-source-hash: fnv1a64:a15c8579415a53a5
% canonical-lessons: RU-C442-ekskursiya, RU-C442-ekskursovod, RU-C442-putevoditel, RU-C442-dostoprimechatelnost, RU-C442-osmatrivat, RU-R442-review-435-442

\chapter{Sightseeing}
\label{ch:ru-sightseeing}
\hlchaptermodality{442}

\begin{chapteropening}
\textbf{By the end of this chapter:} \emph{I can go sightseeing: a guided tour, the guide, the guidebook, the sights and looking round a place.}

Complete the last lesson of chapter 442: I can go sightseeing: a guided tour, the guide, the guidebook, the sights and looking round a place.
\end{chapteropening}

\section[\texorpdfstring{\ru{экскурсия} (ekskúrsiya)}{экскурсия (ekskúrsiya)}]{\ru{экскурсия} (ekskúrsiya) — a guided tour, an excursion}

\label{lesson:RU-C442-ekskursiya}



\begin{quote}
Before the new one: say the Russian for to search for, then the Russian for to ask again.
\end{quote}

\subsection*{You'll want to know: \ru{экскурсия}}
\textbf{\ru{экскурсия}} (\emph{ekskúrsiya}) — \textquotedblleft{}a guided tour, an excursion\textquotedblright{}.

\begin{quote}
\ru{В} \ru{субботу} \ru{у} \ru{нас} \ru{экскурсия} \ru{по} \ru{городу}.
\end{quote}

\begin{quote}
\emph{V subbótu u nas ekskúrsiya po górodu.}
\end{quote}

\begin{quote}
On Saturday we have a guided tour of the town.
\end{quote}

Feminine. \emph{\ru{Экскурсия} \ru{по} \ru{музею}} — a tour of the museum.

\begin{cousinweb}
\emph{\ru{экскурсия}} goes back to Latin \emph{excursiō} 'a running out, a setting out' (EtymDB), from \emph{excurrere} 'to run out'.
\end{cousinweb}

\begin{grammarlens}[title={What you've built}]
One more word for telling, explaining and describing at B1.
\end{grammarlens}

\subsection*{Your turn}
\begin{itemize}
\raggedright
  \item \emph{Say it:} \emph{ekskúrsiya}
  \item \emph{Say it:} \emph{ekskúrsiya}, once more
  \item \emph{Say it:} a sentence of your own with \emph{ekskúrsiya}
  \item \emph{Recall:} say the Russian for to search for, then the Russian for to ask again, then say \emph{ekskúrsiya} again
\end{itemize}

\begin{tcolorbox}[breakable,colback=teal!4,colframe=teal!35!black,title={Before you move on}]
What does \ru{экскурсия} mean? (\textquotedblleft{}A guided tour, an excursion\textquotedblright{}.) Say it once more.
\end{tcolorbox}
\section[\texorpdfstring{\ru{экскурсовод} (ekskursovód)}{экскурсовод (ekskursovód)}]{\ru{экскурсовод} (ekskursovód) — a tour guide}

\label{lesson:RU-C442-ekskursovod}



\begin{quote}
Before the new one: say the Russian for to ask again, then the Russian for a guided tour.
\end{quote}

\subsection*{You'll want to know: \ru{экскурсовод}}
\textbf{\ru{экскурсовод}} (\emph{ekskursovód}) — \textquotedblleft{}a tour guide\textquotedblright{}.

\begin{quote}
\ru{Экскурсовод} \ru{рассказывал} \ru{о} \ru{старых} \ru{домах}.
\end{quote}

\begin{quote}
\emph{Ekskursovód rasskázyval o stárykh domákh.}
\end{quote}

\begin{quote}
The guide talked about the old houses.
\end{quote}

Masculine, for a woman too. Also \emph{\ru{гид}}.

\begin{cousinweb}
\emph{\ru{экскурсовод}} is built from \emph{\ru{экскурсия}} 'a tour' and \emph{-\ru{вод}}, from \emph{\ru{водить}} 'to lead'. \emph{\ru{Экскурсия}} goes back to Latin \emph{excursiō} and \emph{\ru{водить}} to Proto-Slavic \emph{voditi} (EtymDB).
\end{cousinweb}

\begin{grammarlens}[title={What you've built}]
One more word for telling, explaining and describing at B1.
\end{grammarlens}

\subsection*{Your turn}
\begin{itemize}
\raggedright
  \item \emph{Say it:} \emph{ekskursovód}
  \item \emph{Say it:} \emph{ekskursovód}, once more
  \item \emph{Say it:} a sentence of your own with \emph{ekskursovód}
  \item \emph{Recall:} say the Russian for to ask again, then the Russian for a guided tour, then say \emph{ekskursovód} again
\end{itemize}

\begin{tcolorbox}[breakable,colback=teal!4,colframe=teal!35!black,title={Before you move on}]
What does \ru{экскурсовод} mean? (\textquotedblleft{}A tour guide\textquotedblright{}.) Say it once more.
\end{tcolorbox}
\section[\texorpdfstring{\ru{путеводитель} (putevodítel)}{путеводитель (putevodítel)}]{\ru{путеводитель} (putevodítel) — a guidebook}

\label{lesson:RU-C442-putevoditel}



\begin{quote}
Before the new one: say the Russian for a guided tour, then the Russian for a tour guide.
\end{quote}

\subsection*{You'll want to know: \ru{путеводитель}}
\textbf{\ru{путеводитель}} (\emph{putevodítel}) — \textquotedblleft{}a guidebook\textquotedblright{}.

\begin{quote}
\ru{В} \ru{путеводителе} \ru{есть} \ru{хорошая} \ru{карта} \ru{центра}.
\end{quote}

\begin{quote}
\emph{V putevodítele yest khoróshaya kárta tséntra.}
\end{quote}

\begin{quote}
There is a good map of the centre in the guidebook.
\end{quote}

Masculine. Literally a way-leader: \emph{\ru{путь}} + \emph{\ru{водитель}}.

\begin{cousinweb}
\emph{\ru{путеводитель}} is built from \emph{\ru{путь}} 'a way' and \emph{\ru{водитель}} 'one who leads', from \emph{\ru{водить}} 'to lead' (EtymDB): a book that leads you on your way.
\end{cousinweb}

\begin{grammarlens}[title={What you've built}]
One more word for telling, explaining and describing at B1.
\end{grammarlens}

\subsection*{Your turn}
\begin{itemize}
\raggedright
  \item \emph{Say it:} \emph{putevodítel}
  \item \emph{Say it:} \emph{putevodítel}, once more
  \item \emph{Say it:} a sentence of your own with \emph{putevodítel}
  \item \emph{Recall:} say the Russian for a guided tour, then the Russian for a tour guide, then say \emph{putevodítel} again
\end{itemize}

\begin{tcolorbox}[breakable,colback=teal!4,colframe=teal!35!black,title={Before you move on}]
What does \ru{путеводитель} mean? (\textquotedblleft{}A guidebook\textquotedblright{}.) Say it once more.
\end{tcolorbox}
\section[\texorpdfstring{\ru{достопримечательность} (dostoprimechátelnost)}{достопримечательность (dostoprimechátelnost)}]{\ru{достопримечательность} (dostoprimechátelnost) — a sight, a place of interest}

\label{lesson:RU-C442-dostoprimechatelnost}



\begin{quote}
Before the new one: say the Russian for a tour guide, then the Russian for a guidebook.
\end{quote}

\subsection*{You'll want to know: \ru{достопримечательность}}
\textbf{\ru{достопримечательность}} (\emph{dostoprimechátelnost}) — \textquotedblleft{}a sight, a place of interest\textquotedblright{}.

\begin{quote}
\ru{Главная} \ru{достопримечательность} \ru{города} — \ru{старый} \ru{мост}.
\end{quote}

\begin{quote}
\emph{Glávnaya dostoprimechátelnost góroda — stáryy most.}
\end{quote}

\begin{quote}
The main sight of the town is the old bridge.
\end{quote}

Feminine; often plural: \emph{\ru{достопримечательности} \ru{города}} — the sights of the town.

\begin{cousinweb}
\emph{\ru{достопримечательность}} is built from \emph{\ru{досто}-} 'worthy', as in \emph{\ru{достойный}}, and \emph{\ru{примечательный}} 'notable', from \emph{\ru{примечать}} 'to notice': a thing worth noticing. EtymDB does not list it.
\end{cousinweb}

\begin{grammarlens}[title={What you've built}]
One more word for telling, explaining and describing at B1.
\end{grammarlens}

\subsection*{Your turn}
\begin{itemize}
\raggedright
  \item \emph{Say it:} \emph{dostoprimechátelnost}
  \item \emph{Say it:} \emph{dostoprimechátelnost}, once more
  \item \emph{Say it:} a sentence of your own with \emph{dostoprimechátelnost}
  \item \emph{Recall:} say the Russian for a tour guide, then the Russian for a guidebook, then say \emph{dostoprimechátelnost} again
\end{itemize}

\begin{tcolorbox}[breakable,colback=teal!4,colframe=teal!35!black,title={Before you move on}]
What does \ru{достопримечательность} mean? (\textquotedblleft{}A sight, a place of interest\textquotedblright{}.) Say it once more.
\end{tcolorbox}
\section[\texorpdfstring{\ru{осматривать} (osmátrivat)}{осматривать (osmátrivat)}]{\ru{осматривать} (osmátrivat) — to look round (a place), to go sightseeing; to examine}

\label{lesson:RU-C442-osmatrivat}



\begin{quote}
Before the new one: say the Russian for a guidebook, then the Russian for a sight.
\end{quote}

\subsection*{You'll want to know: \ru{осматривать}}
\textbf{\ru{осматривать}} (\emph{osmátrivat}) — \textquotedblleft{}to look round (a place), to go sightseeing; to examine\textquotedblright{}.

\begin{quote}
\ru{Весь} \ru{день} \ru{туристы} \ru{осматривали} \ru{старый} \ru{город}.
\end{quote}

\begin{quote}
\emph{Ves den turísty osmátrivali stáryy górod.}
\end{quote}

\begin{quote}
All day the tourists looked round the old town.
\end{quote}

Imperfective; perfective \emph{\ru{осмотреть}} (\emph{\ru{осмотрю}, \ru{осмотришь}}). Also of a doctor: to examine.

\begin{cousinweb}
\emph{\ru{осматривать}} is built from \emph{\ru{о}-} 'round' and \emph{-\ru{сматривать}}, from \emph{\ru{смотреть}} 'to look', which EtymDB relates to Proto-Slavic \emph{motriti} 'to watch'.
\end{cousinweb}

\begin{grammarlens}[title={What you've built}]
One more word for telling, explaining and describing at B1.
\end{grammarlens}

\subsection*{Your turn}
\begin{itemize}
\raggedright
  \item \emph{Say it:} \emph{osmátrivat}
  \item \emph{Say it:} \emph{osmátrivat}, once more
  \item \emph{Say it:} a sentence of your own with \emph{osmátrivat}
  \item \emph{Recall:} say the Russian for a guidebook, then the Russian for a sight, then say \emph{osmátrivat} again
\end{itemize}

\begin{tcolorbox}[breakable,colback=teal!4,colframe=teal!35!black,title={Before you move on}]
What does \ru{осматривать} mean? (\textquotedblleft{}To look round (a place), to go sightseeing; to examine\textquotedblright{}.) Say it once more.
\end{tcolorbox}
\section[(dialogue)]{Words from chapters 435–442 — a review}

\label{lesson:RU-R442-review-435-442}



\begin{quote}
Say the Russian for service and to serve.
\end{quote}

\subsection*{Your turn}
Say these aloud:

\begin{itemize}
\raggedright
  \item service, to serve, staff, a tip, to be rude
  \item a refund, to demand, to be indignant, to fix, a receipt
  \item to get lost, to find one's way, to roam, a compass, a satnav
  \item a signpost, a sign, an arrow, a landmark, a route map
  \item a turn, a fork, a dead end, to take a short cut, to go round
  \item a crossing, to cross, an embankment, a boulevard, an avenue
  \item a passer-by, to ask around, to tell, to search for, to ask again
  \item a guided tour, a tour guide, a guidebook, a sight, to look round
\end{itemize}

\begin{tcolorbox}[breakable,colback=teal!4,colframe=teal!35!black,title={Before you move on}]
Service? (\textbf{\ru{обслуживание}}.) To serve? (\textbf{\ru{обслуживать}}.) Staff? (\textbf{\ru{персонал}}.) A tip? (\textbf{\ru{чаевые}}.) To be rude? (\textbf{\ru{хамить}}.) A refund? (\textbf{\ru{возврат}}.) To demand? (\textbf{\ru{требовать}}.) To be indignant? (\textbf{\ru{возмущаться}}.) To fix? (\textbf{\ru{устранять}}.) A receipt? (\textbf{\ru{квитанция}}.) To get lost? (\textbf{\ru{заблудиться}}.) To find one's way? (\textbf{\ru{ориентироваться}}.) To roam? (\textbf{\ru{блуждать}}.) A compass? (\textbf{\ru{компас}}.) A satnav? (\textbf{\ru{навигатор}}.) A signpost? (\textbf{\ru{указатель}}.) A sign? (\textbf{\ru{табличка}}.) An arrow? (\textbf{\ru{стрелка}}.) A landmark? (\textbf{\ru{ориентир}}.) A route map? (\textbf{\ru{схема}}.) A turn? (\textbf{\ru{поворот}}.) A fork? (\textbf{\ru{развилка}}.) A dead end? (\textbf{\ru{тупик}}.) To take a short cut? (\textbf{\ru{срезать}}.) To go round? (\textbf{\ru{огибать}}.) A crossing? (\textbf{\ru{переход}}.) To cross? (\textbf{\ru{переходить}}.) An embankment? (\textbf{\ru{набережная}}.) A boulevard? (\textbf{\ru{бульвар}}.) An avenue? (\textbf{\ru{аллея}}.) A passer-by? (\textbf{\ru{прохожий}}.) To ask around? (\textbf{\ru{расспрашивать}}.) To tell? (\textbf{\ru{подсказывать}}.) To search for? (\textbf{\ru{разыскивать}}.) To ask again? (\textbf{\ru{переспрашивать}}.) A guided tour? (\textbf{\ru{экскурсия}}.) A tour guide? (\textbf{\ru{экскурсовод}}.) A guidebook? (\textbf{\ru{путеводитель}}.) A sight? (\textbf{\ru{достопримечательность}}.) To look round? (\textbf{\ru{осматривать}}.)
\end{tcolorbox}
